% =============================================================================
%  preamble.tex -- per-thesis configuration.
%
%  The house style lives in aml-thesis.cls: typography, colours, headings,
%  captions, table and figure conventions, bibliography style. Do not edit that
%  file to write your thesis -- edit this one.
%
%  This file is for the four things that are yours rather than the lab's:
%    1. which .bib file you use
%    2. extra packages your topic needs
%    3. your own notation macros
%    4. deliberate overrides of a class default
%
%  Chapters stay prose-only: a \usepackage or \newcommand in a text/ file
%  belongs here instead.
% =============================================================================

% ---- 1. Bibliography resource -----------------------------------------------
% The citation style is fixed by the class; this is just the file. Add more
% \addbibresource lines if you keep references split across several files.
\addbibresource{bibliography/references.bib}

% ---- 2. Additional packages -------------------------------------------------
% Anything your topic needs and the class does not already load. Check
% aml-thesis.cls first: amsmath, newpxmath, tikz, booktabs, siunitx, listings,
% cleveref, csquotes, subcaption, enumitem, etoolbox and biblatex are all loaded
% there already, and loading a package twice with different options is an error.
% Note that amssymb is deliberately absent: newpxmath supplies the AMS symbols
% and the two clash.
%
% Frequently wanted, deliberately left out of the class:
%   \usepackage[ruled,vlined]{algorithm2e}  % pseudo-code
%   \usepackage{pdfpages}                   % include a signed declaration PDF
%   \usepackage{rotating}                   % landscape (sidewaystable)
%   \usepackage{glossaries}                 % acronyms and a glossary
%
% Load cleveref-sensitive packages with care: cleveref is already loaded by the
% class, and anything that redefines referencing must come after it.

% ---- 3. Your macros ---------------------------------------------------------
% Notation shortcuts. Defining them once here, rather than typing the markup
% out each time, is what lets you change your mind about notation in chapter 4
% without editing chapters 1 through 3.
%
% Examples -- delete or replace:
\newcommand{\vect}[1]{\bm{#1}}              % a vector:  \vect{x}
\newcommand{\mat}[1]{\bm{#1}}               % a matrix:  \mat{X}
\newcommand{\set}[1]{\mathcal{#1}}          % a set:     \set{H}
\newcommand{\reals}{\mathbb{R}}             % the reals: \reals^{d}
\newcommand{\expect}{\mathbb{E}}            % expectation
\DeclareMathOperator*{\argmin}{arg\,min}
\DeclareMathOperator*{\argmax}{arg\,max}

% ---- 4. Overrides -----------------------------------------------------------
% The class defines Paul Tol's seven `vibrant' data colours, and derives the
% document accents from them (accent = tolblue, highlight = tolorange, plus the
% two soft tints). Redefining an accent rebrands the document; redefining the
% tol* colours changes what the figures use. Nothing else hard-codes a colour.
%   \colorlet{accent}{tolteal}
%   \definecolor{tolblue}{HTML}{1F77B4}
%
% Extra theorem environments join the four the class defines (definition,
% theorem, proposition, lemma):
%   \newtheorem{corollary}{Corollary}[chapter]
%
% To write in another language, reload babel here with your language.
